\begin{recipe}{Kaneelrolletjes}
\kicker[Bakken,Snel]{Bakken · snel}
\index[register]{Kaneelrolletjes}
\index[register]{Kaneel}
\index[register]{Croissantdeeg}

\meta{25 minuten \dvd 7 stuks \dvd Oven 180\,°C}

\lettrine{M}{et} kant-en-klaar croissantdeeg staan deze kaneelrolletjes zo op
tafel. Geen rijstijd nodig, alleen uitrollen, insmeren, oprollen en snijden.

\blockrule{Ingrediënten}
\begin{ingredients}
  \ing{2 tl}{kaneel}\index[register]{Kaneel}
  \ing{1 blikje}{croissantdeeg}\index[register]{Croissantdeeg}
  \ing{50 g}{zachte boter}
  \ing{50--60 g}{lichtbruine basterdsuiker}
\end{ingredients}

\medskip
\noindent\textit{Glazuur}
\begin{ingredients}
  \ing{100 g}{poedersuiker}
  \ing{1 tl}{vanille-extract}
  \ing{2 el}{melk}
\end{ingredients}

\blockrule{Bereiding}
\tip{Rol het deeg uit tot één vierkant in plaats van de driehoeken waar het
blikje voor bedoeld is, dan krijg je een mooie, gelijkmatige rol.}
\begin{steps}
  \item Verwarm de oven voor op ca.\ 180\,°C (boven- en onderwarmte), of
        volgens de temperatuur op de verpakking van het croissantdeeg.
  \item Rol het croissantdeeg uit tot één vierkant, druk de naden dicht.
  \item Mix de zachte boter met de kaneel en de basterdsuiker tot een
        smeerbare pasta.
  \item Smeer de kaneelboter gelijkmatig over het deeg en rol het vierkant
        strak op tot een rol.
  \item Snijd de rol in gelijke plakken.
  \item Leg de rolletjes met wat afstand van elkaar in een ingevette
        ovenschaal, zodat ze tijdens het bakken kunnen uitzetten.
  \item Bak 15--20 minuten tot de rolletjes goudbruin zijn.
  \item Klop ondertussen de poedersuiker, het vanille-extract en de melk tot
        een glad glazuur. Verdeel het glazuur over de rolletjes zodra ze uit
        de oven komen.
\end{steps}
\end{recipe}
